\section{September 23, 2026}

\subsection{Injectivity, surjectivity, and dimension}

Recall from Proposition~\ref{prop:injective-kernel} that a linear map
$T\in\mathcal L(V,W)$ is injective if and only if $\ker T=\{0\}$.
In particular, if $\ker T=\{0\}$, then
\[
  T(v)=T(v')
  \quad\Longrightarrow\quad T(v-v')=0
  \quad\Longrightarrow\quad v-v'=0
  \quad\Longrightarrow\quad v=v'.
\]
The corresponding criterion for surjectivity follows directly from
the definition of the image:
\[
  T\text{ is surjective}
  \quad\Longleftrightarrow\quad \operatorname{im}T=W.
\]

\subsection{Proof of the fundamental theorem of linear maps}

We now prove Theorem~\ref{thm:fundamental-linear-maps}: if $V$ is
finite-dimensional and $T\in\mathcal L(V,W)$, then
$\operatorname{im}T$ is finite-dimensional and
\[
  \dim V=\dim\ker T+\dim\operatorname{im}T.
\]
This identity is also called the \emph{rank--nullity formula};
$\dim\operatorname{im}T$ is the rank and $\dim\ker T$ is the nullity.

\begin{proof}[Proof of Theorem~\ref{thm:fundamental-linear-maps}]
  Since $\ker T$ is a subspace of the finite-dimensional space $V$,
  it has a basis $u_1,\ldots,u_m$. Extend this collection to a basis
  \[
    u_1,\ldots,u_m,v_1,\ldots,v_n
  \]
  of $V$. Thus $\dim\ker T=m$ and $\dim V=m+n$. We claim that
  \[
    T(v_1),\ldots,T(v_n)
  \]
  is a basis of $\operatorname{im}T$.

  First, every $v\in V$ can be written as
  \[
    v=a_1u_1+\cdots+a_mu_m+b_1v_1+\cdots+b_nv_n.
  \]
  Because each $u_i$ belongs to $\ker T$, linearity gives
  \[
    T(v)=b_1T(v_1)+\cdots+b_nT(v_n).
  \]
  Hence $\operatorname{im}T$ is contained in the span of the displayed
  images. The reverse inclusion holds because each $T(v_j)$ belongs
  to the subspace $\operatorname{im}T$. Therefore
  \[
    \operatorname{im}T
      =\operatorname{span}\{T(v_1),\ldots,T(v_n)\}.
  \]

  To show linear independence, suppose that
  \[
    c_1T(v_1)+\cdots+c_nT(v_n)=0.
  \]
  Then $T(c_1v_1+\cdots+c_nv_n)=0$, so
  $c_1v_1+\cdots+c_nv_n\in\ker T$. There are scalars
  $d_1,\ldots,d_m$ such that
  \[
    c_1v_1+\cdots+c_nv_n=d_1u_1+\cdots+d_mu_m.
  \]
  Rearranging gives
  \[
    -d_1u_1-\cdots-d_mu_m+c_1v_1+\cdots+c_nv_n=0.
  \]
  The chosen basis of $V$ is linearly independent, so all the
  coefficients are zero. In particular, $c_1=\cdots=c_n=0$.

  Thus the image has a basis of $n$ vectors, so
  $\dim\operatorname{im}T=n$. It follows that
  \[
    \dim V=m+n=\dim\ker T+\dim\operatorname{im}T.
  \]
  The same argument includes $m=0$ or $n=0$, using the empty basis
  for the zero subspace.
\end{proof}

\begin{figure}[htbp]
  \centering
  \begin{tikzpicture}[>=Stealth, font=\small,
      box/.style={draw, rounded corners, minimum width=3.6cm,
        minimum height=1cm, align=center}]
    \node[box, fill=red!7] (kernel) at (0,0)
      {$u_1,\ldots,u_m$\\basis of $\ker T$};
    \node[box, fill=blue!7] (complement) at (0,-1.6)
      {$v_1,\ldots,v_n$\\extension to a basis of $V$};
    \node[box, fill=red!7] (zero) at (6,0) {$0$};
    \node[box, fill=blue!7] (range) at (6,-1.6)
      {$T(v_1),\ldots,T(v_n)$\\basis of $\operatorname{im}T$};
    \node at (0,1) {$V$};
    \node at (6,1) {$W$};
    \draw[->, thick] (kernel) -- node[above] {$T$} (zero);
    \draw[->, thick] (complement) -- node[above] {$T$} (range);
  \end{tikzpicture}
  \caption{The kernel basis maps to zero; the images of the remaining
    basis vectors form a basis of the image.}
  \label{fig:rank-nullity-bases}
\end{figure}

\subsection{Dimension restrictions on linear maps}

\begin{corollary}\label{cor:dimension-obstructs-injectivity}
  Let $V$ and $W$ be finite-dimensional vector spaces. If
  $\dim V>\dim W$, then no linear map from $V$ to $W$ is injective.
\end{corollary}

\begin{proof}
  If $T\in\mathcal L(V,W)$, then $\operatorname{im}T\subseteq W$, so
  $\dim\operatorname{im}T\leq\dim W$. Rank--nullity gives
  \[
    \dim\ker T
      =\dim V-\dim\operatorname{im}T
      \geq\dim V-\dim W>0.
  \]
  Hence $\ker T\neq\{0\}$, and $T$ is not injective.
\end{proof}

\begin{example}
  Define $T\colon\bb{R}^4\to\bb{R}^3$ by
  \[
    T\begin{bmatrix}x_1\\x_2\\x_3\\x_4\end{bmatrix}
      =\begin{bmatrix}
        \sqrt{7}\,x_1+\pi x_2+x_4\\
        97x_1+3x_2+2x_3\\
        x_2+6x_3+7x_4
      \end{bmatrix}.
  \]
  Since $4>3$, this map is not injective, regardless of the particular
  coefficients. There is no need to solve for its kernel explicitly.
  The same example works over $\bb{C}$.
\end{example}

\begin{proposition}\label{prop:dimension-obstructs-surjectivity}
  Suppose that $V$ and $W$ are finite-dimensional and
  $\dim V<\dim W$. Then no linear map from $V$ to $W$ is surjective.
\end{proposition}

\begin{proof}
  For $T\in\mathcal L(V,W)$, rank--nullity gives
  \[
    \dim\operatorname{im}T
      =\dim V-\dim\ker T
      \leq\dim V<\dim W.
  \]
  Thus $\operatorname{im}T\neq W$, so $T$ is not surjective.
\end{proof}

\subsection{A complement to the kernel}

\begin{proposition}\label{prop:kernel-complement-image}
  Let $V$ be finite-dimensional, and let $T\in\mathcal L(V,W)$.
  There is a subspace $U\subseteq V$ such that
  \[
    V=\ker T\oplus U
    \qquad\text{and}\qquad
    \operatorname{im}T=\{T(u):u\in U\}.
  \]
  In particular, $U\cap\ker T=\{0\}$, and the restriction
  $T|_U\colon U\to\operatorname{im}T$ is bijective.
\end{proposition}

\begin{proof}
  As in the proof of rank--nullity, extend a basis $u_1,\ldots,u_m$
  of $\ker T$ to a basis $u_1,\ldots,u_m,v_1,\ldots,v_n$ of $V$.
  Set
  \[
    U=\operatorname{span}\{v_1,\ldots,v_n\}.
  \]
  The combined basis spans $V$, so $V=\ker T+U$. If $w\in U\cap\ker T$,
  then
  \[
    w=b_1v_1+\cdots+b_nv_n=a_1u_1+\cdots+a_mu_m.
  \]
  Linear independence of the combined basis forces every coefficient
  to vanish. Thus $w=0$, and the sum is direct.

  Every $v\in V$ has the form $v=k+u$, where $k\in\ker T$ and
  $u\in U$. Therefore $T(v)=T(u)$, proving that $T|_U$ maps onto
  $\operatorname{im}T$. Its kernel is $U\cap\ker T=\{0\}$, so it
  is also injective.
\end{proof}

\begin{figure}[htbp]
  \centering
  \begin{tikzpicture}[>=Stealth, font=\small, scale=1.05]
    \draw[->, thick, blue!65!black] (-1,0) -- (3.5,0)
      node[right] {$U$};
    \draw[->, thick, red!70!black] (0,-.6) -- (0,2.8)
      node[above] {$\ker T$};
    \draw[dashed, gray] (2.5,0) -- (2.5,1.8) -- (0,1.8);
    \draw[->, thick] (0,0) -- (2.5,1.8)
      node[above right] {$v=(a,b)$};
    \draw[->, thick, blue!65!black] (0,0) -- (2.5,0);
    \draw[->, thick, red!70!black] (0,0) -- (0,1.8);
    \fill (2.5,0) circle (1.5pt);
    \node[below] at (2.5,0) {$u=(a,0)$};
    \node[left] at (0,1.8) {$k=(0,b)$};
    \node[below left] at (0,0) {$0$};
  \end{tikzpicture}
  \caption{For $T\colon\bb{R}^2\to\bb{R}$, $T(x,y)=x$, take $U$ to
    be the horizontal axis. The decomposition $v=k+u$ gives
    $T(v)=T(u)=a$, while $T(k)=0$.}
  \label{fig:kernel-complement-projection}
\end{figure}
